\section{Planning: Searching for the Best Plan}
% The "how to search" methods + MPC. Comes AFTER the cliff arc: the mechanism
% (how-to-plan.tex) was already shown before the cliff; here we make the search
% practical and finally name the replan loop as MPC.

\begin{frame}{Back to Planning: How to Search?}
\begin{itemize}
    \item Recall: planning = pick the action sequence whose model roll-out scores
          highest. The hard part was the $\arg\max$ over sequences
    \item The fix we just used (replan + collect data) also relied on this search
    \item Two separate questions: first, how to \emph{search} for a plan (random
          shooting, CEM); then, how often to \emph{replan} (MPC)
\end{itemize}
\end{frame}

% === Random shooting ====================================================
\begin{frame}{How to Search? Random Shooting}
\begin{itemize}
    \item Simplest idea (no gradients, just guess):
    \begin{enumerate}
        \item Sample $N$ action sequences $a_{t:t+H}^{(1)}, \dots, a_{t:t+H}^{(N)}$
              from some distribution
        \item Roll each one through the model $f_\phi$, sum the predicted rewards
        \item Keep the sequence with the highest predicted return
    \end{enumerate}
\end{itemize}
\end{frame}

\begin{frame}{How to Search? Random Shooting}
\begin{itemize}
    \item Simplest idea (no gradients, just guess):
    \begin{enumerate}
        \item Sample $N$ action sequences $a_{t:t+H}^{(1)}, \dots, a_{t:t+H}^{(N)}$
              from some distribution
        \item Roll each one through the model $f_\phi$, sum the predicted rewards
        \item Keep the sequence with the highest predicted return
    \end{enumerate}
    \item Derivative-free and trivially parallel
    \item This is the planner in \textbf{Nagabandi et al.}: ``random shooting'',
          no ensembles
    \item Weakness: most samples land in low-reward regions, wasteful
\end{itemize}
\end{frame}

% === CEM ================================================================
\begin{frame}{Cross-Entropy Method (CEM)}
\begin{itemize}
    \item Fix random shooting's waste: \emph{refit} the sampling distribution
          toward the good sequences
\end{itemize}
\end{frame}

\begin{frame}{Cross-Entropy Method (CEM)}
\begin{itemize}
    \item Fix random shooting's waste: \emph{refit} the sampling distribution
          toward the good sequences
    \begin{enumerate}
        \item Sample $N$ sequences from a Gaussian $\mathcal{N}(\mu, \Sigma)$
        \item Evaluate under the model; keep the top-$k$ \emph{elites}
        \item Refit $\mu, \Sigma$ to the elites
        \item Repeat for a few iterations
    \end{enumerate}
    \item The distribution concentrates on high-return plans
    \item This is the optimizer hiding inside step 3 of our algorithm
\end{itemize}
\end{frame}

% === MPC: a different question, NOT a search method =====================
\section{Planning: When to Replan?}
\begin{frame}{A Different Question: When to Replan?}
\begin{itemize}
    \item Random shooting and CEM answered \emph{how to find} a good plan
    \item MPC answers a \emph{different} question: not how to \emph{search},
          but \textbf{how much of a plan to trust before re-deciding}
    \item It is a planning \emph{technique}, not another search method
\end{itemize}
\end{frame}

% === MPC ================================================================
\begin{frame}{Model Predictive Control (MPC)}
\begin{itemize}
    \item A full open-loop plan is fragile: model errors compound over the
          horizon $H$
\end{itemize}
\end{frame}

\begin{frame}{Model Predictive Control (MPC)}
\definecolor{mbgreen}{RGB}{122,173,75}
\definecolor{mborange}{RGB}{237,148,72}
\definecolor{mbblue}{RGB}{79,129,189}
\begin{itemize}
    \item A full open-loop plan is fragile: model errors compound over the
          horizon $H$
    \item \textbf{MPC}: plan a horizon, but \emph{execute only the first action},
          then re-plan from the state you actually land in
\end{itemize}
\vspace{0.5em}
\centering
\begin{tikzpicture}[>={Stealth[length=2.4mm]}, font=\scriptsize]
    \foreach \i in {0,1,2,3}{
        \fill[mbblue] (1.7*\i,0) circle (2.4pt);
    }
    \node[below] at (0,-0.15) {$s_t$};
    \node[below] at (1.7,-0.15) {$s_{t+1}$};
    \node[below] at (3.4,-0.15) {$s_{t+2}$};
    \node[below] at (5.1,-0.15) {$s_{t+3}$};
    % planned (faded) horizon from s_t
    \draw[mborange!55, dashed, ->] (0,0) .. controls (0.8,0.9) and (1.4,0.9) .. (1.7,0);
    \draw[mborange!35, dashed] (1.7,0) .. controls (2.5,0.8) .. (3.4,0.2);
    \draw[mborange!25, dashed] (3.4,0.2) .. controls (4.2,0.8) .. (5.1,0.1);
    % executed (solid) first step
    \draw[mbgreen, very thick, ->] (0,0) .. controls (0.8,-0.7) and (1.4,-0.7) .. (1.7,0);
    \node[mbgreen!60!black] at (0.85,-0.95) {execute 1st action};
    \node[mborange!70!black] at (3.4,1.0) {planned horizon (discarded)};
\end{tikzpicture}
\end{frame}

\begin{frame}{Model Predictive Control (MPC)}
\begin{itemize}
    \item A full open-loop plan is fragile: model errors compound over the
          horizon $H$
    \item \textbf{MPC}: plan a horizon, but \emph{execute only the first action},
          then re-plan from the state you actually land in
    \item This \emph{is} the ``execute, then re-plan'' loop we used to fix the
          cliff, now it has a name
    \item Next: combine a CEM planner with an \emph{ensemble} model
          $\Rightarrow$ \textbf{PETS}
\end{itemize}
\end{frame}
